% Generated by generate_method_documentation.py.

\subsection{Caso ilustrativo: explicação sobre TUSD}

\label{sec:caso_ilustrativo}

O caso \texttt{CASE-015} foi selecionado por ter a menor minuta final ao consumidor entre os 34 bloqueios julgados inadequados. O critério favorece uma apresentação compacta; o exemplo não representa a frequência dos motivos no conjunto.

\paragraph*{Pergunta.} O que é essa tal de TUSD na fatura? É taxa extra?

\paragraph*{Recuperação.} O Método 3 recuperou os seguintes trechos, na ordem registrada:

\begin{enumerate}

\item Leitura de fatura > Anatomia da fatura para o diagnóstico comercial > 4. Componentes tarifários: TE e TUSD, identificador \path{pub-61a26859e8808af2999a}, similaridade 0,773.

\item Distribuidoras e tarifas de energia > Composição tarifária: TUSD e TE, identificador \path{pub-24c4f657fff094437ead}, similaridade 0,756.

\item Leitura de fatura > O que é descontável vs. O que permanece fixo, identificador \path{pub-956e3a26cfe47d8878fe}, similaridade 0,738.

\end{enumerate}

O primeiro trecho contém a evidência:

\begin{quote}Custo do transporte e manutenção da rede física.\end{quote}

\paragraph*{Tentativa inicial.} A classificação por unidades foi aprovada. A auditoria factual foi rejeitada pelo validador com o código \texttt{citation\_not\_found}. O fluxo efetuou uma reformulação.

\paragraph*{Minuta final ao consumidor.}

\begin{quote}A TUSD não é uma taxa extra! Ela significa Tarifa de Uso do Sistema de Distribuição e faz parte da fatura regular. De forma simples: a TE remunera a energia consumida, e a TUSD cobre o custo de transporte e manutenção da rede física até o seu imóvel.\end{quote}

\paragraph*{Segunda revisão.} A classificação por unidades foi aprovada novamente. A auditoria factual foi rejeitada com \texttt{claim\_not\_in\_draft}, porque uma alegação devolvida pelo auditor não era um trecho literal da unidade correspondente. O contraste preservado no registro é:

\begin{description}

\item[Unidade fornecida:] low\_construal\_template: A TUSD não é uma cobrança extra ou penalidade! A sigla significa Tarifa de Uso do Sistema de Distribuição. A conta de luz homologada pela ANEEL é dividida essencialmente em duas partes: a TE (Tarifa de Energia), que remunera a energia elétrica consumida, e a TUSD, que cobre o transporte dessa energia e a manutenção da rede física (postes e cabos). Deseja que eu analise sua fatura para detalhar esses itens, ou prefere que um consultor humano o auxilie nessa conferência?

\item[Alegação extraída:] A conta de luz homologada pela ANEEL é dividida essencialmente em duas partes: a TE (Tarifa de Energia)... e a TUSD

\end{description}

O validador exige que a alegação extraída seja uma substring exata da unidade. Essa verificação é distinta da conferência de citações nas fontes, que admite normalização. A rejeição estrutural do parecer não equivale a demonstrar erro factual na minuta.

\paragraph*{Decisão e avaliação humana.} O sistema bloqueou a entrega e registrou recomendação de atendimento humano. A latência total foi de 108,59 s. O avaliador marcou o bloqueio como incorreto, indicando que a minuta era globalmente adequada e poderia ser entregue. O exemplo distingue a validação do parecer automático da avaliação global do texto; não permite estender o mesmo motivo aos outros bloqueios.

Os três trechos completos, as duas orientações estruturadas, os pareceres e o julgamento humano estão preservados no arquivo \path{caso-ilustrativo.json}, em \path{output/revisao/documentacao-metodo-2026-09-29/}.
